%% ma: L10-B14-0051
%% lop: 10
%% chuong: 5
%% bai: 14
%% dang: 10.14.9
%% loai: TL
%% muc: TH
%% dap_an: "-"
%% nguon: "chuyen-de-cac-so-dac-trung-cua-mau-so-lieu-khong-ghep-nhom-toan-10-chuong-trinh-moi.pdf (GV. Phan Nhật Linh) (P05), Bài 03 Dạng 2 Bài tập tự luận 11"
%% kien_thuc: "Lớp 10 Bài 13 (số trung bình), Bài 14 (phương sai, độ lệch chuẩn)"
%% trang_thai: chinh-thuc
%% ghi_chu: "Tự duyệt theo yêu cầu giáo viên 22/10/2026."
\begin{ex}
Trong $5$ lần nhảy cao, hai bạn Hà và Mai có kết quả (đơn vị: cm) như sau:
\begin{center}
\begin{tabular}{|c|c|c|c|c|c|}\hline
Hà & $85$ & $90$ & $85$ & $100$ & $90$\\\hline
Mai & $85$ & $90$ & $90$ & $90$ & $95$\\\hline
\end{tabular}
\end{center}
Kết quả nhảy cao của bạn nào ổn định hơn?
\loigiai{$\overline x_H=\overline x_M=90$.

$s_H^2=\dfrac{25+0+25+100+0}5=30$; $s_M^2=\dfrac{25+0+0+0+25}5=10$.

Vì $s_M^2<s_H^2$ nên kết quả của Mai ổn định hơn.}
\end{ex}
